% TEXTO DO ARTIGO
% ================
% Autores: escrevam o conteúdo nas seções abaixo e mantenham a ordem. A abertura
% é criada a partir de metadados.tex; não repita título, autores ou resumo aqui.
% METADADOS DO ARTIGO
% ===================
% Este é o único arquivo em que autores preenchem dados editoriais. Substitua
% todos os campos por informações verificadas. Os dados abaixo são FICTÍCIOS e
% existem apenas para que este modelo possa ser visualizado antes do envio.

\newcommand{\ArtigoTitulo}{Lorem ipsum: demonstração de estrutura para artigos científicos}
\newcommand{\ArtigoSubtitulo}{Texto ilustrativo para revisão editorial e diagramação}

% Liste autores na ordem de autoria e use sobrescritos para associá-los às
% filiações. Mantenha asterisco somente para a pessoa correspondente.
\newcommand{\ArtigoAutores}{Autora de Exemplo\textsuperscript{1*}; Coautor de Exemplo\textsuperscript{2}}
\newcommand{\ArtigoAfiliacoes}{\textsuperscript{1}Instituição de Exemplo, Campus Demonstrativo, Cidade, BA, Brasil.\\\textsuperscript{2}Instituição de Exemplo, Campus Ilustrativo, Cidade, BA, Brasil.}
\newcommand{\ArtigoContatos}{E-mail: autora.exemplo@exemplo.org. ORCID: https://orcid.org/0000-0000-0000-0000.}

% Resumo em português: máximo de 250 palavras; deve informar objetivo e os
% principais resultados. Palavras-chave: até cinco, sem repetir o título.
\newcommand{\ArtigoResumo}{Este texto fictício demonstra a organização prevista para a submissão de artigos no livro. Seu objetivo é permitir que autores e equipe de diagramação observem a abertura, as seções, as citações, as figuras e as referências antes da substituição pelo manuscrito definitivo. Não apresenta resultados de pesquisa nem informações institucionais reais.}
\newcommand{\ArtigoPalavrasChave}{Diagramação. Tipografia. Manuscrito. Demonstração.}

% Traduções fiéis do título, resumo e palavras-chave em português.
\newcommand{\ArtigoTituloIngles}{Lorem ipsum: a structural demonstration for scientific articles}
\newcommand{\ArtigoAbstract}{This fictional text demonstrates the organisation required for articles in the book. It lets authors and the layout team inspect the opening matter, sections, citations, figures, and references before replacing it with the final manuscript. It reports neither research results nor real institutional information.}
\newcommand{\ArtigoKeywords}{Typesetting. Typography. Manuscript. Demonstration.}

\aberturaartigo

\section{Introdução}
% APAGUE este texto demonstrativo e escreva o seu manuscrito nesta estrutura.
Jogos digitais combinam programas, dispositivos de processamento e formas de
interação. Enquanto a partida acontece, o sistema recebe comandos do jogador,
atualiza o estado do jogo e apresenta respostas na tela e nos alto-falantes.
Este exemplo aborda, de maneira introdutória, dois componentes relacionados a
esse ciclo: a unidade central de processamento (CPU) e o controle de jogo.

A Figura~\ref{fig:cpu} representa a CPU como componente de processamento. Em
um jogo, o processador executa partes da lógica e coordena tarefas do sistema;
as funções exatas variam conforme o jogo e a plataforma. A citação
\parencite{exemplo2026} é fictícia e serve somente para demonstrar a chamada
bibliográfica no texto.

\inserirfigura[0.36\textwidth]
  {fig:cpu}
  {CPU como componente de processamento de um jogo digital}
  {imagens/imagem_01.png}
  {Imagem ilustrativa do modelo (2026).}

\section{Fundamentação teórica}
% Desenvolva conceitos e trabalhos que sustentam a pesquisa.
O processamento de um jogo envolve tarefas como interpretar regras,
atualizar personagens e coordenar a comunicação com os demais componentes.
Em computadores e consoles, essas responsabilidades são distribuídas entre
unidades e subsistemas diferentes. Por isso, a imagem da CPU é uma ilustração
conceitual, não um diagrama da arquitetura de um equipamento específico.

O controle, por sua vez, transforma ações físicas — como pressionar um botão
ou mover um analógico — em sinais que o jogo pode interpretar. A Tabela~\ref{tab:elementos-modelo}
resume esses papéis no ciclo básico de interação.

\begin{table}[H]
  \centering
  \caption{\label{tab:elementos-modelo}Componentes do ciclo de interação em um jogo digital}
  \begin{tabular}{p{0.30\textwidth} p{0.64\textwidth}}
    \toprule
    \textbf{Componente} & \textbf{Papel ilustrativo} \\
    \midrule
    Controle & Registra os comandos enviados pelo jogador. \\
    CPU & Processa partes da lógica e coordena tarefas do jogo. \\
    Software do jogo & Interpreta os comandos e atualiza o estado da partida. \\
    \bottomrule
  \end{tabular}
  \legend{Fonte: Síntese demonstrativa elaborada para o modelo (2026).}
\end{table}

\section{Material e métodos}
% Descreva desenho, dados, ferramentas, procedimentos e recorte do estudo.
Este texto é uma demonstração editorial, não um estudo experimental. Para
montá-lo, foram observados os dois ícones incluídos na pasta de imagens e
organizadas suas funções em uma tabela conceitual. Nenhum jogo, console ou
controle foi testado, e não foram coletados dados de desempenho ou de uso.

A Figura~\ref{fig:controle} ilustra um controle acompanhado por um circuito
eletrônico. Ela ajuda a mostrar como uma figura pode ser mencionada no texto,
receber título e fonte e permanecer próxima ao trecho relacionado.

\inserirfigura[0.42\textwidth]
  {fig:controle}
  {Controle de jogo e circuito eletrônico ilustrativos}
  {imagens/imagem_02.png}
  {Imagem ilustrativa do modelo (2026).}

\section{Resultados e discussão}
% Apresente resultados verificáveis e discuta-os à luz da literatura.
As imagens e a tabela destacam um fluxo geral: o jogador envia uma entrada pelo
controle, o software interpreta o comando e o sistema atualiza a partida. A
CPU participa do processamento, embora o trabalho real seja compartilhado com
outros componentes. Como o exemplo usa somente ícones, ele não permite comparar
desempenho, medir latência ou avaliar a experiência de jogadores.

\section{Considerações finais}
% Responda ao objetivo sem introduzir resultados novos.
O exemplo mostra como apresentar CPU e controle como partes do ciclo de
interação em jogos digitais. As figuras ilustram os componentes e a tabela
resume seus papéis. O conteúdo é apenas demonstrativo: um artigo de pesquisa
deve substituir essas descrições por fontes, métodos e resultados próprios.

\section*{Contribuição dos autores}
% Substitua por papéis CRediT reais de cada autor.
Autora de Exemplo: conceituação, redação e visualização. Coautor de Exemplo:
revisão e supervisão.

\section*{Conflito de interesses}
Os autores declaram não haver conflito de interesses.

\section*{Aprovação ética}
Não se aplica.

\referenciasartigo
% Garante que a última página do artigo seja composta ainda com o cabeçalho
% corrente antes de a inclusão desta pasta terminar. Não remova esta linha.
\thispagestyle{fancy}
\clearpage
